% Generated by scripts/locobench_paper_examples.py from f003.
\paragraph{Question.} How did astronomers establish whether the ``Great Debate'' nebulae lay within the Milky Way or formed separate galaxies, and what evidence ultimately settled the dispute? (Astronomy, Natural Sciences.)

\paragraph{Claims.} The plan's claims in assertion order, with their factuality labels; claims joined by a ladder's slots are boxed.

\begin{center}\scriptsize
\setlength{\tabcolsep}{2.5pt}
\begin{tabular}{@{}rp{0.76\linewidth}l@{}}
\toprule
& Claim & Label \\
\midrule
$a_{1}$ & The Great Debate was held on 26 April 1920 at the Smithsonian Museum of Natural History in Washington, D.C. & true \\
$a_{2}$ & Harlow Shapley argued that the spiral nebulae were gaseous clouds lying inside a single large Milky Way system. & true \\
$a_{3}$ & Heber Curtis argued that the spiral nebulae were distant stellar systems comparable in size to the Milky Way. & true \\
\slotclaim{$a_{4}$} & Curtis cited novae observed in spiral nebulae as evidence that those nebulae lay at great distances. & true \\
\slotclaim{$a_{5}$} & Curtis's analysis of novae in spiral nebulae placed M31 about 10,000 light years from the Sun. & \textbf{false} \\
$a_{6}$ & S Andromedae, which appeared in 1885, was a supernova far more luminous than an ordinary nova. & true \\
$a_{7}$ & Shapley relied on Adriaan van Maanen's reported measurements of internal rotation in spiral nebulae. & true \\
$a_{8}$ & Van Maanen's measurements of rotation in spiral nebulae were confirmed by later photographic plate comparisons. & \textbf{false} \\
$a_{9}$ & A later reexamination of van Maanen's plates concluded that the reported rotations of spiral nebulae were measurement artifacts. & true \\
$a_{10}$ & Henrietta Leavitt established the period-luminosity relation from Cepheid variables in the Magellanic Clouds. & true \\
\slotclaim{$a_{11}$} & Edwin Hubble identified Cepheid variable stars in M31 using the 100-inch Hooker telescope at Mount Wilson. & true \\
\slotclaim{$a_{12}$} & Hubble derived a distance to M31 of roughly 900,000 light years. & true \\
$a_{13}$ & M31 lies outside the boundaries of the Milky Way system. & true \\
$a_{14}$ & M31 lies within the boundaries of the Milky Way system. & true \\
\bottomrule
\end{tabular}
\end{center}

\paragraph{Template.} The answer as written once, with its slot markers.

\begin{quote}\small
The status of the spiral nebulae was first contested in open argument and only afterwards resolved by measurement. The Great Debate was held on 26 April 1920 at the Smithsonian Museum of Natural History in Washington, D.C. Harlow Shapley argued that the spiral nebulae were gaseous clouds lying inside a single large Milky Way system, whereas Heber Curtis argued that the spiral nebulae were distant stellar systems comparable in size to the Milky Way. The issue therefore turned on distance, and distance could be extracted only from objects of known intrinsic brightness.

\slotmarker{S1} By contrast, S Andromedae, which appeared in 1885, was a supernova far more luminous than an ordinary nova. The two classes of outburst were not yet separated in the photometry then available, so the eruptions recorded in the spirals could not yet be made to yield a secure scale.

Shapley's case for a single, all-encompassing system rested on evidence of a quite different kind, namely apparent motion within the spirals themselves, since a nebula visibly turning on a human timescale would have to be small and close at hand. \slotmarker{S2} A later reexamination of van Maanen's plates concluded that the reported rotations of spiral nebulae were measurement artifacts, and that verdict cannot be squared with any claim that subsequent comparisons of photographic plates had confirmed those rotations.

Resolution came from a calibrated distance indicator rather than from debate. Earlier, Henrietta Leavitt established the period-luminosity relation from Cepheid variables in the Magellanic Clouds; the search for such stars in the spirals came only afterwards. \slotmarker{S3} A figure of that magnitude demonstrates that M31 lies outside the boundaries of the Milky Way system. The two placements exclude one another: either M31 lies beyond the Galaxy's confines, or M31 lies within the boundaries of the Milky Way system, and exactly one of these must hold. On the Cepheid evidence, the spirals are stellar systems in their own right.
\end{quote}

\paragraph{Fillers.} Each slot's fillers, as admitted by the committee.

\begin{center}\scriptsize
\setlength{\tabcolsep}{2.5pt}
\begin{tabular}{@{}lp{0.80\linewidth}@{}}
\toprule
Filler & Text \\
\midrule
\multicolumn{2}{@{}l}{\slotmarker{S1} sign slot, $a_{4}$--$a_{5}$} \\
\quad conflict & Curtis cited novae observed in spiral nebulae as evidence that those nebulae lay at great distances; that evidence indicates a remove for M31 of about 10,000 light years from the Sun, the figure his analysis of those novae yielded. \\
\quad no link & Curtis cited novae observed in spiral nebulae as evidence that those nebulae lay at great distances. His analysis of novae in spiral nebulae placed M31 about 10,000 light years from the Sun. \\
\quad support & Curtis cited novae observed in spiral nebulae as evidence that those nebulae lay at great distances, and that evidence rules out the claim, sometimes attached to his analysis of novae in spiral nebulae, that it placed M31 about 10,000 light years from the Sun. \\
\addlinespace
\multicolumn{2}{@{}l}{\slotmarker{S2} sign (spare) slot, $a_{7}$--$a_{8}$} \\
\quad conflict & Shapley relied on Adriaan van Maanen's reported measurements of internal rotation in spiral nebulae, and that reliance shows that the rotations van Maanen reported were confirmed by later comparisons of photographic plates. \\
\quad no link & Shapley relied on Adriaan van Maanen's reported measurements of internal rotation in spiral nebulae. The rotations van Maanen reported were confirmed by later comparisons of photographic plates. \\
\quad support & Shapley relied on Adriaan van Maanen's reported measurements of internal rotation in spiral nebulae, and because nothing beyond those original reports of rotation stood behind the figures, the rotations cannot have been confirmed by later comparisons of photographic plates. \\
\addlinespace
\multicolumn{2}{@{}l}{\slotmarker{S3} control slot, $a_{11}$--$a_{12}$} \\
\quad wording 1 & Edwin Hubble identified Cepheid variable stars in M31 using the 100-inch Hooker telescope at Mount Wilson, and that identification led him to a distance for the nebula of roughly 900,000 light years. \\
\quad wording 2 & Because Edwin Hubble was able to pick out Cepheid variable stars in M31 with the 100-inch Hooker telescope at Mount Wilson, a distance of roughly 900,000 light years followed for that nebula. \\
\quad wording 3 & A distance to M31 of roughly 900,000 light years was the direct result of Edwin Hubble's detection of Cepheid variable stars in that nebula with the 100-inch Hooker telescope at Mount Wilson. \\
\bottomrule
\end{tabular}
\end{center}

\paragraph{Rungs.} The five responses of the family, which filler each carries in the edited slots, and gold $\LCS$ at the verdict and the asserted world.

\begin{center}\scriptsize
\setlength{\tabcolsep}{2.5pt}
\begin{tabular}{@{}l>{\raggedright\arraybackslash}p{0.27\linewidth}>{\raggedright\arraybackslash}p{0.30\linewidth}cc@{}}
\toprule
& Ladder, rung & Fillers & \multicolumn{2}{c}{$\LCS$ at world} \\
& & & verdict & asserted \\
\midrule
t0 & sign $\rho_0$ & S1: conflict & 0.123 & 0.036 \\
t1 & sign $\rho_1$, control $\rho_0$ & S1: no link & 0.292 & 0.030 \\
t2 & sign $\rho_2$ & S1: support & 0.334 & 0.036 \\
t3 & control $\rho_1$ & S3: wording 2 & 0.292 & 0.030 \\
t4 & control $\rho_2$ & S3: wording 3 & 0.292 & 0.030 \\
\bottomrule
\end{tabular}
\end{center}

\paragraph{Response t1 in full} (sign $\rho_1$; edited slots highlighted, 406 words).

\begin{quote}\small
The status of the spiral nebulae was first contested in open argument and only afterwards resolved by measurement. The Great Debate was held on 26 April 1920 at the Smithsonian Museum of Natural History in Washington, D.C. Harlow Shapley argued that the spiral nebulae were gaseous clouds lying inside a single large Milky Way system, whereas Heber Curtis argued that the spiral nebulae were distant stellar systems comparable in size to the Milky Way. The issue therefore turned on distance, and distance could be extracted only from objects of known intrinsic brightness.

\slotfill{S1}{Curtis cited novae observed in spiral nebulae as evidence that those nebulae lay at great distances. His analysis of novae in spiral nebulae placed M31 about 10,000 light years from the Sun.} By contrast, S Andromedae, which appeared in 1885, was a supernova far more luminous than an ordinary nova. The two classes of outburst were not yet separated in the photometry then available, so the eruptions recorded in the spirals could not yet be made to yield a secure scale.

Shapley's case for a single, all-encompassing system rested on evidence of a quite different kind, namely apparent motion within the spirals themselves, since a nebula visibly turning on a human timescale would have to be small and close at hand. Shapley relied on Adriaan van Maanen's reported measurements of internal rotation in spiral nebulae. The rotations van Maanen reported were confirmed by later comparisons of photographic plates. A later reexamination of van Maanen's plates concluded that the reported rotations of spiral nebulae were measurement artifacts, and that verdict cannot be squared with any claim that subsequent comparisons of photographic plates had confirmed those rotations.

Resolution came from a calibrated distance indicator rather than from debate. Earlier, Henrietta Leavitt established the period-luminosity relation from Cepheid variables in the Magellanic Clouds; the search for such stars in the spirals came only afterwards. \slotfill{S3}{Edwin Hubble identified Cepheid variable stars in M31 using the 100-inch Hooker telescope at Mount Wilson, and that identification led him to a distance for the nebula of roughly 900,000 light years.} A figure of that magnitude demonstrates that M31 lies outside the boundaries of the Milky Way system. The two placements exclude one another: either M31 lies beyond the Galaxy's confines, or M31 lies within the boundaries of the Milky Way system, and exactly one of these must hold. On the Cepheid evidence, the spirals are stellar systems in their own right.
\end{quote}

\paragraph{Gold relations of response t1.} Sense, compiled coupling, strength band, and the relation's sign at the verdict (v.) and asserted (a.) world ($+$ support, $-$ conflict, $\cdot$ no factor); a slot id marks a relation drawn by a filler.

\begin{center}\scriptsize
\setlength{\tabcolsep}{2.2pt}
\begin{tabular}{@{}lllllcc@{}}
\toprule
Pair & & Sense & Coupling & Band & v. & a. \\
\midrule
$a_{2}\to a_{3}$ &  & Contrast & none & strong & $\cdot$ & $\cdot$ \\
$a_{4}\to a_{6}$ &  & Contrast & none & moderate & $\cdot$ & $\cdot$ \\
$a_{7}\to a_{2}$ &  & Evidence & entailment & moderate & $+$ & $+$ \\
$a_{8}\to a_{9}$ &  & Refutation & contradiction & strong & $+$ & $-$ \\
$a_{10}\to a_{11}$ &  & Precedence & none & moderate & $\cdot$ & $\cdot$ \\
$a_{11}\to a_{12}$ & S3 & Cause-Effect & entailment & strong & $+$ & $+$ \\
$a_{12}\to a_{13}$ &  & Evidence & entailment & strong & $+$ & $+$ \\
$a_{13}\to a_{14}$ &  & Alternative & exclusive & strong & $-$ & $-$ \\
\bottomrule
\end{tabular}
\end{center}
